\documentclass{article}

\usepackage{amsmath}
\usepackage{booktabs}
\usepackage{tabularx}

\title{Supplementary Material Outline: BodyFutures-Bench}
\author{Anonymous Authors}

\begin{document}
\maketitle

\section{Environment Specifications}

The supplement will provide the complete equations and locked configuration
values for each BodyFutures environment.

\begin{table}[h]
\centering
\caption{Supplement environment equation checklist.}
\begin{tabularx}{\linewidth}{lX}
\toprule
Environment & Required details \\
\midrule
FeverAnt & Temperature, energy, torque damage, actuator weakening, FAC,
collapse thresholds, observation modes. \\
WoundWalker & Wound accumulation, impact load, balance health, energy,
actuator weakening, FAC, collapse thresholds. \\
ThermalHopper & Heat from action and landing impact, cooling, damage,
energy, actuator weakening, FAC. \\
NumbArm & Numbness, proprioceptive corruption, manipulation load,
damage, energy, action rescaling, success criterion, FAC. \\
\bottomrule
\end{tabularx}
\end{table}

\section{Preregistered Protocol}

This section will reproduce the frozen CoRL manifest:
seeds, train budgets, episode lengths, evaluation horizons, baseline list,
checkpoint-selection score, IWM dataset mixture, coverage gates, and final
metric definitions.  Any deviation from the manifest will be listed with a
version identifier and reason.

\section{Run Commands}

The supplement will include command-line recipes for reproducing the final
experiments:
\begin{verbatim}
bash experiments/run_corl_2026_rescue_primary.sh
bash experiments/run_corl_2026_primary.sh
bash experiments/run_corl_2026_secondary.sh
bash experiments/run_corl_2026_all.sh
\end{verbatim}
For archival final sweeps, the exact manifest version, output root, run id,
and seed set will be reported.

\section{Weights and Biases Metric Schema}

All scalar metrics will be logged against a shared global step.  The supplement
will list stable metric keys for training, rollout diagnostics, body state,
evaluation, inference, IWM training, and policy video.  Step-specific video
filenames may exist locally for traceability, but W\&B keys will remain stable.

\section{IWM Dataset Coverage}

For each environment and seed, the supplement will report:
\begin{itemize}
    \item transition count and episode count by collection mode;
    \item FAC minimum, mean, standard deviation, and histogram;
    \item fraction of samples below FAC thresholds $0.75$ and $0.50$;
    \item collapse-adjacent fraction and collapse labels;
    \item whether automatic coverage repair was invoked.
\end{itemize}

\section{IWM Architecture and Training}

The supplement will provide model dimensions, ensemble size, input layout,
horizon heads, loss weights, optimizer settings, batch size, device placement,
and calibration diagnostics.  For final runs, ensemble size and training
settings will be frozen before seed sweeps.

\section{Evaluation Details}

The supplement will define:
\begin{itemize}
    \item task reward per step and cumulative task return;
    \item survival fraction and collapse rate;
    \item final FAC, minimum FAC, viability cost, and body score;
    \item early-warning time before collapse;
    \item prediction RMSE, low-FAC error, calibration, AUROC, and AUPRC;
    \item bootstrap or Student confidence interval computation.
\end{itemize}

\section{Compute Budget}

The final supplement will report CPU model, GPU model, RAM, number of vector
environments, wall-clock runtime per method, IWM training time, and total compute
used for final runs.  Pilot and debugging compute will be reported separately
from frozen result compute.

\section{Video Storyboard}

The video supplement will show:
\begin{enumerate}
    \item FeverAnt task PPO sprinting early and degrading over long horizons;
    \item FeverAnt PPO-IWM preserving FAC while retaining task progress;
    \item WoundWalker comparison under long-horizon degradation;
    \item IWM calibration visualization and FAC trajectories;
    \item representative failure cases where IWM is overconservative or
    miscalibrated.
\end{enumerate}

\section{Failure Cases and Negative Results}

The supplement will explicitly include runs where PPO-IWM fails to improve over
current-viability PPO or loses too much task reward.  These cases are important
for identifying whether the contribution is stronger as a benchmark and
prediction protocol, or as an algorithm that reliably improves the full
task-viability Pareto frontier.

\end{document}
